\documentclass[../../main.tex]{subfiles}
\begin{document}
\section{流形上的辛结构}\label{section:ch22-symplectic-structures-on-manifolds}

现在转向光滑流形 $M$。$M$ 上的一个非退化 $2$-形式是这样的 $2$-形式 $\omega$：对每个 $p\in M$，$\omega_p$ 都是非退化的 $2$-余向量。$M$ 上的一个\emph{辛形式}是一个闭的非退化 $2$-形式。赋予了一个特定辛形式选取的光滑流形称为\emph{辛流形}。辛形式的选取有时也称为一个\emph{辛结构}。

\cref{proposition:ch22-symplectic-canonical-form}表明辛流形必须是偶维的。然而，并非所有偶维光滑流形都容许辛结构。例如，\cref{proposition:ch22-omega-nonzero}表明：若 $\omega$ 是 $2n$ 流形上的辛形式，则 $\omega^n$ 是一个处处非零的 $2n$-形式，故每个辛流形都是可定向的。此外，习题 22-5 中描述了一个必要的同调条件。它特别地蕴含：$S^2$ 是唯一容许辛结构的球面。

若 $(M_1,\omega_1)$ 与 $(M_2,\omega_2)$ 是辛流形，满足 $F^*\omega_2=\omega_1$ 的微分同胚 $F:M_1\to M_2$ 称为\emph{辛微分同胚}（symplectomorphism）。研究辛流形在辛微分同胚下不变的性质，称为\emph{辛几何}或\emph{辛拓扑}。

\begin{example}[辛流形]\label{example:ch22-symplectic-manifolds}
\begin{enumerate}[(a)]
\item 在 $\mathbb{R}^{2n}$ 上取标准坐标记作 $(x_1,\dots,x_n,y_1,\dots,y_n)$，$2$-形式
\begin{align}
\omega=\sum_{i=1}^n \mathrm{d}x^i\wedge\mathrm{d}y^i
\end{align}
是辛的：它显然是闭的；它是非退化的，因为它在每一点处的值就是\cref{example:ch22-standard-symplectic-tensor}中的辛张量。这称为 $\mathbb{R}^{2n}$ 上的\emph{标准辛形式}。（在涉及标准辛形式的公式中，就像涉及 Euclid 内积的公式一样，通常需要显式写出求和号，因为求和指标 $i$ 在上位置出现了两次。）
\item 设 $\Sigma$ 是任意可定向光滑 $2$ 流形，$\omega$ 是 $\Sigma$ 上处处非零的光滑 $2$-形式。则 $\omega$ 是闭的，因为 $\mathrm{d}\omega$ 是 $2$ 流形上的 $3$-形式。此外，正如我们在\cref{example:ch22-standard-symplectic-tensor}之后刚观察到的，在二维中每个处处非零的 $2$-形式都是非退化的，故 $(\Sigma,\omega)$ 是辛流形。
\end{enumerate}
\end{example}

设 $(M,\omega)$ 是辛流形。一个（浸入的或嵌入的）子流形 $S\subset M$，如果 $T_pS$（看作 $T_pM$ 的子空间）在每一点 $p\in S$ 都具有相应性质，就称为辛的、迷向的、余迷向的或 Lagrang 子流形。更一般地，一个光滑浸入（或嵌入）$F:N\to M$，若子空间 $\mathrm{d}F_p(T_pN)\subset T_{F(p)}M$ 对每个 $p\in N$ 都具有相应性质，就称它具有这些性质之一。因此，一个子流形是辛的（迷向的等），当且仅当其包含映射具有相同性质。

\begin{exercise}\label{exercise:ch22-10}
设 $(M,\omega)$ 是辛流形，$F:N\to M$ 是光滑浸入。证明：$F$ 是迷向的当且仅当 $F^*\omega=0$；$F$ 是辛的当且仅当 $F^*\omega$ 是辛形式。
\end{exercise}

\subsection*{余切从上的典范辛形式}

最重要的辛流形是余切丛的全空间，它们带有我们现在要定义的典范辛结构。首先，$T^*M$ 的全空间上有一个自然的 $1$-形式 $\sigma$，称为\emph{ tautological $1$-form}（重言 $1$-形式），定义如下。$T^*M$ 中的一点是某个 $q\in M$ 处的一个余向量 $\varphi\in T^*_qM$；我们把这样一点记作 $(q,\varphi)$。自然投影 $\pi:T^*M\to M$ 就是 $\pi(q,\varphi)=q$，它在 $q$ 处的逐点拉回是线性映射 $\mathrm{d}\pi^*_{(q,\varphi)}:T^*_qM\to T^*_{(q,\varphi)}(T^*M)$。我们定义 $\sigma\in\Omega^1(T^*M)$（$T^*M$ 全空间上的 $1$-形式）为
\begin{align}
\sigma_{(q,\varphi)}=\mathrm{d}\pi^*_{(q,\varphi)}\varphi.
\end{align}
换言之，$\sigma$ 在 $(q,\varphi)\in T^*M$ 处的值就是余向量 $\varphi$ 本身关于 $\pi$ 的拉回。若 $v$ 是 $T_{(q,\varphi)}(T^*M)$ 中的切向量，则
\begin{align*}
\sigma_{(q,\varphi)}(v)=\varphi\big(\mathrm{d}\pi_{(q,\varphi)}(v)\big).
\end{align*}

\begin{figure}[ht]
\centering
\includegraphics[width=0.8\textwidth]{fig-22-1.png}
\caption{$T^*M$ 上的重言 1-形式}\label{figure:ch22-tautological}
\end{figure}

\begin{proposition}\label{proposition:ch22-canonical-symplectic}
设 $M$ 是光滑流形。重言 $1$-形式 $\sigma$ 是光滑的，而 $\omega=-\mathrm{d}\sigma$ 是 $T^*M$ 全空间上的辛形式。
\end{proposition}

\begin{proof}
设 $(x^i)$ 是 $M$ 上的光滑坐标，并令 $(x^i,\xi_i)$ 表示 $T^*M$ 上相应的自然坐标（如第 277 页所定义）。回忆：$(q,\varphi)\in T^*M$ 的坐标定义为 $(x^i,\xi_i)$，其中 $(x^i)$ 是 $q$ 的坐标表示，而 $\xi_i\,\mathrm{d}x^i$ 是 $\varphi$ 的坐标表示。在这些坐标下，投影 $\pi:T^*M\to M$ 的坐标表达式为 $\pi(x,\xi)=x$。这意味着 $\mathrm{d}\pi^*(\mathrm{d}x^i)=\mathrm{d}x^i$，故 $\sigma$ 的坐标表达式为
\begin{align}
\sigma_{(x,\xi)}=\mathrm{d}\pi^*_{(x,\xi)}(\xi_i\,\mathrm{d}x^i)=\xi_i\,\mathrm{d}x^i.
\end{align}
由此立刻可知 $\sigma$ 是光滑的，因为它在这些坐标下的分量函数是线性的。

令 $\omega=-\mathrm{d}\sigma\in\Omega^2(T^*M)$。显然 $\omega$ 是闭的，因为它是恰当的。此外，在自然坐标下，由上式得
\begin{align*}
\omega=\sum_i \mathrm{d}x^i\wedge\mathrm{d}\xi_i.
\end{align*}
在借助这些坐标把 $T^*M$ 的一个开子集与 $\mathbb{R}^{2n}$ 的一个开子集等同起来之后，$\omega$ 对应于 $\mathbb{R}^{2n}$ 上的标准辛形式（用 $\xi_i$ 代替 $y_i$）。由此可知 $\omega$ 是辛的。
\end{proof}

这个命题中定义的辛形式称为 $T^*M$ 上的\emph{典范辛形式}。它的众多用途之一，是给出“一个 $1$-形式为闭”这一说法的一个略更“几何化”的解释。

\begin{proposition}\label{proposition:ch22-closed-1form-lagrangian}
设 $M$ 是光滑流形，$\beta$ 是 $M$ 上的光滑 $1$-形式。把 $\beta$ 看作从 $M$ 到 $T^*M$ 的光滑映射，则 $\beta$ 是光滑嵌入，而 $\beta$ 是闭的当且仅当其像 $\beta(M)$ 是 $T^*M$ 的一个 Lagrang 子流形。
\end{proposition}

\begin{proof}
在整个证明中，我们需要记住 $\beta:M\to T^*M$ 同时扮演两个角色：一方面它是 $M$ 上的 $1$-形式；另一方面它是流形之间的光滑映射。由于它们实际上是同一个映射，我们不使用不同记号来区分；但你应小心地在论证的每一步都想清楚 $\beta$ 在扮演哪个角色。

用 $M$ 上的光滑局部坐标 $(x^i)$ 以及 $T^*M$ 上相应的自然坐标 $(x^i,\xi_i)$，映射 $\beta:M\to T^*M$ 的坐标表示为
\begin{align*}
\beta(x^1,\dots,x^n)=\big(x^1,\dots,x^n,\beta_1(x),\dots,\beta_n(x)\big),
\end{align*}
其中 $\beta_i\,\mathrm{d}x^i$ 是 $\beta$ 作为 $1$-形式的坐标表示。由此立刻可知 $\beta$ 是光滑浸入；它是单射的，因为 $\pi\circ\beta=\operatorname{Id}_M$。要证明它是嵌入，由命题 4.22 只需证明它是固有映射。而这又由 $\pi$ 是 $\beta$ 的左逆这一事实、借助命题 A.53 得到。

由于 $\beta(M)$ 是 $n$ 维的，它是 Lagrang 的当且仅当它是迷向的，而这当且仅当 $\beta^*\omega=0$。$\beta$ 对重言形式 $\sigma$ 的拉回为
\begin{align*}
\beta^*\sigma=\beta^*(\xi_i\,\mathrm{d}x^i)=\beta_i\,\mathrm{d}x^i=\beta.
\end{align*}
这一点也可以从下面这个略更不变的计算中看出：
\begin{align*}
(\beta^*\sigma)_p(v)=\sigma_{\beta(p)}(\mathrm{d}\beta_p(v))
=\beta_p(\mathrm{d}\pi_{\beta(p)}\circ\mathrm{d}\beta_p(v))=\beta_p(v),
\end{align*}
这由 $\sigma$ 的定义以及 $\pi\circ\beta=\operatorname{Id}_M$ 得到。因此，
\begin{align*}
\beta^*\omega=\beta^*(-\mathrm{d}\sigma)=-\mathrm{d}(\beta^*\sigma)=-\mathrm{d}\beta.
\end{align*}
由此可知 $\beta$ 是 Lagrang 嵌入当且仅当 $\mathrm{d}\beta=0$。
\end{proof}
\end{document}
